\section{Range projections and their defects}
\label{sec:perturbations_projections}

For a closed subspace $Z$ of a Hilbert space, $P_Z$ denotes the orthogonal
projection onto $Z$, and $P_T$ the orthogonal projection onto the range of a
map $T$.

\begin{lemma}[Continuity of injective range projectors]
    \label{thm:continuous_injective_range_projector}
    \lean{ContinuousLinearMap.continuous_of_euclideanSingle,
        ContinuousLinearMap.continuousAt_range_starProjection_of_injective,
        ContinuousLinearMap.continuous_injectiveRangeProjector}
    \leanok
    Let $E$ be a finite-dimensional Hilbert space, $F$ a Hilbert space, and
    $T_x\colon E\to F$ a family of bounded maps continuous at $x_0$, with
    $T_{x_0}$ injective. Then $x\mapsto P_{T_x}$ is continuous at $x_0$; no
    injectivity of $T_x$ for $x\neq x_0$ is assumed. In particular, for a
    continuous family of injective maps, the orthogonal projection
    $T_x(T_x^\dagger T_x)^{-1}T_x^\dagger$ onto the range of $T_x$ depends
    continuously on $x$. A family of maps out of $\C^I$, for a finite set $I$,
    is continuous when its values on the standard basis vectors are.
\end{lemma}
\begin{proof}
    \leanok
    Injectivity is an open condition, so near $x_0$ the projection is
    $T_x(T_x^\dagger T_x)^{-1}T_x^\dagger$. The Gram operator
    $T_x^\dagger T_x$ is continuous in $x$ and invertible at $x_0$, and
    inversion is continuous at invertible elements of a Banach algebra. The
    last statement holds because a vector is a finite combination of the
    standard basis vectors.
\end{proof}

\begin{lemma}[Supported Gram inverses and range projections]
    \label{thm:supported_gram_projector_cancellation}
    \lean{ContinuousLinearMap.comp_supportedGramInverse_adjoint_eq_starProjection,
        ContinuousLinearMap.starProjection_comp_sub_of_supportedGramFactorizations}
    \leanok
    Let $T:E\to\mathcal H$ be a continuous linear map from a finite-dimensional
    Hilbert space to a complete Hilbert space. Suppose $Q=Q^*$, $TQ=T$, and
    $(T^*T)H=Q$. Then $P_T=THT^*$.

    More generally, let $T:E\to\mathcal H$, $L:F\to\mathcal H$, and
    $C:G\to\mathcal H$ have finite-dimensional domains and a common
    factorization $C=TJ_1=LJ_2$. For $S\in\{T,L,C\}$, suppose
    $Q_S=Q_S^*$, $SQ_S=S$, and $(S^*S)H_S=Q_S$. If also $H_T=H_T^*$, then
    \begin{align}
        P_TP_L-P_C
         & =TH_T\bigl[T^*L-(T^*T)J_1H_CJ_2^*(L^*L)\bigr]H_LL^*.
        \label{eq:ph_supported_gram_residual}
    \end{align}
    This is the algebra behind the projector comparison in
    \cite[Lemma~6.1(ii)]{Nachtergaele1996Spectral}.
\end{lemma}
\begin{proof}\leanok
    The vector $THT^*x$ lies in $\operatorname{ran}T$. Taking adjoints of
    $TQ=T$ gives $QT^*=T^*$, so
    $T^*(x-THT^*x)=T^*x-QT^*x=0$. Thus $THT^*$ is the orthogonal range
    projection. For the three-map identity, self-adjointness of $H_T$ gives
    $H_T(T^*T)=Q_T$. Consequently
    $TH_T(T^*T)=T$ and $(L^*L)H_LL^*=L^*$.
    Expanding the right-hand side of~(\ref{eq:ph_supported_gram_residual})
    and using $TJ_1=C$ and $J_2^*L^*=C^*$ gives
    $P_TP_L-CH_CC^*$.
\end{proof}

\begin{lemma}[Norm control from supported projector cancellation]
    \label{thm:supported_gram_projector_norm}
    \lean{ContinuousLinearMap.norm_starProjection_comp_sub_le_of_supported_factorization,
        ContinuousLinearMap.tendsto_norm_starProjection_comp_sub_zero_of_supported_factorizations}
    \leanok
    Let $T:E\to F$, $L:G\to F$, and $C:W\to F$ be bounded
    maps, with finite-dimensional domains and a complete Hilbert space $F$.
    Suppose
    \begin{align}
        P_TP_L-P_C=T H_T R H_L L^*.
        \notag
    \end{align}
    Then
    \begin{align}
        \|P_TP_L-P_C\|
        \le \|T\|\|H_T\|\|H_L\|\|L\|\|R\|.
        \notag
    \end{align}
    For a family along any filter, assume this factorization eventually
    holds, the product $\|T\|\|H_T\|\|H_L\|\|L\|$ is eventually
    bounded, and $\|R\|\to0$. Then $\|P_TP_L-P_C\|\to0$.
    The spaces may vary with the index.
\end{lemma}
\begin{proof}\leanok
    Apply submultiplicativity of the operator norm and $\|L^*\|=\|L\|$.
    The eventual bound then gives a constant multiple of $\|R\|$;
    squeeze the nonnegative projection norm between zero and this quantity.
\end{proof}

\begin{lemma}[Three-operator perturbation estimate]
    \label{thm:three_operator_perturbation}
    \lean{ContinuousLinearMap.norm_comp_three_sub_le}
    \leanok
    For bounded linear operators $A,a:G\to W$, $B,b:F\to G$, and
    $C,c:E\to F$ between normed spaces,
    \begin{align}
        \|ABC-abc\|\le
        \|A-a\|\|B\|\|C\|+
        \|a\|\|B-b\|\|C\|+
        \|a\|\|b\|\|C-c\|.
    \end{align}
    No finite-dimensional or completeness hypothesis is required.
\end{lemma}
\begin{proof}\leanok
    Use $ABC-abc=(A-a)BC+a(B-b)C+ab(C-c)$, the triangle inequality,
    and submultiplicativity of the operator norm.
\end{proof}

\begin{lemma}[Vanishing three-operator perturbations]
    \label{thm:three_operator_perturbation_limit}
    \lean{ContinuousLinearMap.tendsto_norm_comp_three_sub_zero_of_bounds}
    \leanok
    \uses{thm:three_operator_perturbation}
    Let the normed spaces and operators in Lemma~\ref{thm:three_operator_perturbation}
    vary along an arbitrary filter. Suppose the three individual differences
    tend to zero in norm, and the products $\|B\|\|C\|$, $\|a\|\|C\|$,
    and $\|a\|\|b\|$ have a common eventual upper bound.
    Then $\|ABC-abc\|\to0$.
\end{lemma}
\begin{proof}\leanok
    \uses{thm:three_operator_perturbation}
    Bound the product difference by the common bound times the sum of the
    three individual norm differences, and apply the squeeze theorem.
\end{proof}

\begin{lemma}[Isometric conjugation of a projection defect]
    \label{thm:isometric_projection_defect_transport}
    \lean{Submodule.norm_map_starProjection_comp_sub_eq}
    \leanok
    Let $U,V,W$ be closed subspaces of a Hilbert space and let $F$ be a
    surjective linear isometry to another Hilbert space. Then
    \begin{align}
        \|P_{FU}P_{FV}-P_{FW}\|=\|P_W-P_VP_U\|.
    \end{align}
\end{lemma}
\begin{proof}\leanok
    Orthogonal projections are conjugated by $F$, and conjugation by a
    surjective isometry preserves the operator norm. Taking the adjoint
    reverses the order of the two projections and preserves the norm.
\end{proof}

\begin{lemma}[Finite sector projection estimate]
    \label{thm:finite_sector_projection_estimate}
    \lean{Submodule.norm_starProjection_comp_sub_le_sector_errors}
    \leanok
    Let $U,V,W$ and $S_i,T_i,C_i$ be subspaces of a finite-dimensional
    real or complex Hilbert space, with $i$ ranging over a finite set $I$.
    Set
    $e_U=\|P_U-\sum_iP_{S_i}\|$,
    $e_V=\|P_V-\sum_iP_{T_i}\|$, and
    $e_W=\|\sum_iP_{C_i}-P_W\|$. Then
    \begin{align}
        \|P_UP_V-P_W\|
        \leq e_U+|I|e_V+
        \sum_{i,j\in I}\|P_{S_i}P_{T_j}-\delta_{ij}P_{C_i}\|+e_W.
    \end{align}
    No orthogonality of the sector subspaces is required.
    This finite-sum estimate assembles the projection comparisons in
    \cite[Lemmas~commutation(ii) and disjoint]{Nachtergaele1996Spectral}.
\end{lemma}
\begin{proof}\leanok
    Put $P=\sum_iP_{S_i}$, $Q=\sum_iP_{T_i}$, and $Z=\sum_iP_{C_i}$.
    Expand
    $P_UP_V-P_W=(P_U-P)P_V+P(P_V-Q)+(PQ-Z)+(Z-P_W)$.
    Every orthogonal projection has norm at most one, so $\|P\|\leq|I|$.
    Expand $PQ-Z$ as the stated double sum and apply the triangle inequality.
\end{proof}

\begin{lemma}[Vanishing sector defects give a vanishing joint projection defect]
    \label{thm:sector_projection_defect_limit}
    \lean{Submodule.tendsto_norm_starProjection_comp_sub_zero_of_sector_errors}
    \leanok
    \uses{thm:finite_sector_projection_estimate}
    Let the subspaces in Lemma~\ref{thm:finite_sector_projection_estimate}
    vary along an arbitrary filter, allowing the finite-dimensional ambient
    Hilbert spaces to vary. Suppose the three projection-sum errors tend to
    zero, and suppose $\|P_{S_i}P_{T_i}-P_{C_i}\|\to0$ for each $i$, while
    $\|P_{S_i}P_{T_j}\|\to0$ for each $i\ne j$.
    Then $\|P_UP_V-P_W\|\to0$.
    The finite index set is fixed and may be empty.
\end{lemma}
\begin{proof}\leanok
    \uses{thm:finite_sector_projection_estimate}
    The finite double sum of sector errors tends to zero.
    Apply Lemma~\ref{thm:finite_sector_projection_estimate} and squeeze the
    nonnegative joint defect between zero and its vanishing upper bound.
\end{proof}
